\section{Research Methodology}\label{sec:method}

\subsection{Research Design}

A controlled computational experiment with a factorial design was carried out. The model (four
models), the evaluation design (random split, temporal split, district holdout), the missingness
pattern (three patterns), the missingness rate (0, 10, 20, 30 and 40\%) and the strategy used to
handle missing values (four strategies) were treated as the \emph{independent variables}. The
prediction error on held-out observations, measured by the metrics of Section~\ref{sec:metrics},
was treated as the \emph{dependent variable}. All other factors were held fixed: the
hyperparameters of every model were kept the same in every condition, and within each crop and
seed the same missingness masks were reused for every model and strategy, so that all comparisons
between models and strategies are paired. Feature selection and location features were studied in
two additional experiments with the same protocol.

\subsection{Participants}

No human participants or personal data were involved, so ethical approval was not required. The
units of analysis are district--crop--season--year observations compiled from public, aggregated
government statistics, as described in Section~\ref{sec:data}.

\subsection{Proposed System}

The experimental framework is shown in Fig.~\ref{fig:framework}. District crop statistics and daily
weather are cleaned and joined, and seasonal weather features are computed for each observation.
For each crop, a missingness mask is drawn and applied, the observations are divided into training
and test folds according to the evaluation design, the missing-value strategy and the model are
fitted on the training folds only, and predictions are made for the test fold. The out-of-fold
predictions are pooled, and the metrics are computed and compared across conditions.

\begin{figure}[!t]
\centering
\begin{tikzpicture}[node distance=3.2mm, every node/.style={font=\scriptsize},
  box/.style={draw, rounded corners=1pt, align=center, minimum width=5.2cm, inner sep=2.2pt},
  arr/.style={-{Latex[length=1.6mm]}, thin}]
\node[box] (crop) {DES district crop statistics\\ + NASA POWER daily weather (1997--2023)};
\node[box, below=of crop] (clean) {Cleaning, area filter, join by district and season};
\node[box, below=of clean] (feat) {Seasonal weather features (18 per observation)};
\node[box, below=of feat] (mask) {Missingness mask: MCAR / variable-wise / block outage};
\node[box, below=of mask] (split) {Evaluation design: random / temporal / district holdout};
\node[box, below=of split] (fit) {Missing-value strategy + model fitted on training folds};
\node[box, below=of fit] (pred) {Pooled out-of-fold predictions $\rightarrow$ RMSE, MAE, $R^2$, nRMSE};
\node[box, below=of pred] (stat) {Bootstrap over districts, Friedman test across conditions};
\foreach \a/\b in {crop/clean, clean/feat, feat/mask, mask/split, split/fit, fit/pred, pred/stat}
  \draw[arr] (\a) -- (\b);
\node[draw, dashed, rounded corners=2pt, fit=(mask)(split)(fit), inner sep=2.5pt] (rep) {};
\node[right=1mm of rep, align=left, font=\scriptsize] {repeated\\ for 3 seeds};
\end{tikzpicture}
\caption{Experimental framework used for every crop, model, strategy and condition.}
\label{fig:framework}
\end{figure}

\subsection{Dataset Collection}\label{sec:data}

Crop statistics were obtained from the district-wise, season-wise Area, Production and Yield (APY)
data of the Directorate of Economics and Statistics (DES), Ministry of Agriculture and Farmers
Welfare \cite{desapy}, which cover the agricultural years 1997--98 to 2022--23 and report area in
hectares and production in tonnes. Gujarat was chosen as the study region, and three major crops
with long and consistent records were selected: groundnut and bajra, grown in the Kharif
(monsoon) season, and wheat, grown in the Rabi (winter) season. Daily weather was obtained from the
NASA POWER daily point API \cite{nasapower} for 1~January 1997 to 31~December 2023 at the
headquarters of each of Gujarat's 33 districts. The meteorological fields of NASA POWER are derived
from the MERRA-2 reanalysis on a grid of $0.5^\circ \times 0.625^\circ$, and solar radiation from
the CERES product. Seven variables were retrieved: corrected precipitation, mean, maximum and
minimum air temperature at 2~m, relative humidity at 2~m, all-sky downward shortwave radiation and
wind speed at 2~m. The resulting dataset is summarised in Table~\ref{tab:data}.

\begin{table}[!t]
\caption{Analysis Dataset After Cleaning}\label{tab:data}
\centering\footnotesize\setlength{\tabcolsep}{3.5pt}
\begin{tabular}{llrrcc}
\toprule
Crop & Season (window) & Obs. & Distr. & Years & Yield (t/ha) \\
\midrule
Groundnut & Kharif (1 Jun--31 Oct) & 450 & 29 & 1997--2023 & 1.64 $\pm$ 0.84 \\
Bajra & Kharif (1 Jun--31 Oct) & 420 & 25 & 1997--2023 & 1.28 $\pm$ 0.57 \\
Wheat & Rabi (1 Nov--31 Mar) & 618 & 32 & 1997--2023 & 2.77 $\pm$ 0.69 \\
\bottomrule
\end{tabular}


\vspace{2pt}{\scriptsize Obs.: district--year observations; Distr.: districts; yield as mean $\pm$
standard deviation.}
\end{table}

\subsection{Data Preprocessing}

Of the 455,359 records in the national file, 19,407 belong to Gujarat and 1,888 to the three
selected crop--season combinations (groundnut and bajra in Kharif, wheat in Rabi). One record
with missing production was removed. Records with a cropped area below 1,000~ha were then
excluded, because yields computed from very small areas are unreliable; for example, a groundnut
record of 49~t/ha was found among them. This rule removed 399 records and left 1,488 observations
from 32 districts; Navsari, where none of the three crops reached 1,000~ha, dropped out. No
duplicates were present, and yield was computed as production divided by area. District names in
the crop file and the weather requests were matched exactly, including the seven districts formed
in 2013, which have records only from 2014--15. The raw weather contained no missing values, so all
missing values in the experiments are simulated.

\subsection{Feature Engineering}

For each observation, the daily weather was cut to the season window of its crop: 1~June to
31~October (153 days) for Kharif and 1~November to 31~March (151 days) for Rabi. Eighteen features
were computed, grouped as follows. \emph{Rainfall} (7): seasonal total, number of rainy days
($\geq 2.5$~mm, the threshold of the India Meteorological Department), standard deviation of daily
rainfall, maximum one-day rainfall and totals for the early, middle and late thirds of the season.
\emph{Temperature} (7): seasonal means of the daily mean, maximum and minimum temperature, number of
hot days (maximum above 35\,$^\circ$C in Kharif and 32\,$^\circ$C in Rabi) and mean temperature in
each third of the season. \emph{Humidity}, \emph{solar radiation} and \emph{wind} (1 each): seasonal
means. The \emph{year} of the season was added to capture the technology trend in yields. All
features were computed with missing-aware operations: means and counts used the available days,
totals were rescaled to the full window, and a feature was set to missing when more than half of
its days were missing.

\subsection{Missing-Data Simulation}

Let $\rho \in \{0.1, 0.2, 0.3, 0.4\}$ denote the missingness rate. Three mechanisms were simulated.
\begin{itemize}
\item \emph{Random (MCAR):} each weather feature of each observation is removed independently with
probability $\rho$.
\item \emph{Variable-wise:} for each observation, each of the five variable groups (rainfall,
temperature, humidity, solar radiation, wind) is removed as a whole with probability $\rho$,
which mimics a sensor or data product that is unavailable for a season.
\item \emph{Block outage:} for each observation and each of the seven daily variables, a contiguous
run of $\lfloor \rho L \rceil$ days is removed from the $L$-day season window at a uniformly random
position, and the features are then recomputed from the incomplete daily record.
\end{itemize}
Under the first two mechanisms, features become missing; under the block mechanism, most features
remain available but are estimated from fewer days, and only features of a season third that loses
more than half of its days become missing. The same mechanism and rate were applied to training and
test observations, with independent draws, because historical records used for training are also
incomplete in practice. The year feature was never removed.

\subsection{Missing-Value Strategies}

Four strategies were compared, each fitted on the training folds only: (i)~\emph{median}
replacement; (ii)~\emph{KNN} imputation with $k = 5$ neighbours on standardised features
\cite{troyanskaya2001}; (iii)~\emph{iterative} imputation, in which each feature is regressed on
the others with Bayesian ridge regression for up to ten rounds \cite{vanbuuren2011}; and
(iv)~\emph{native} handling in XGBoost, in which missing values are sent along a default branch
learned during training \cite{chen2016}. Strategies (i)--(iii) were combined with every model and
strategy (iv) with XGBoost only, giving ten model--strategy combinations.

\subsection{Models and Training Configuration}

The four models and their fixed settings are listed in Table~\ref{tab:models}. A mean predictor
serves as the naive reference. Ridge regression \cite{hoerl1970} fits a linear relation with an
$\ell_2$ penalty chosen by internal cross-validation. In a random forest \cite{breiman2001}, many
regression trees are grown on bootstrap samples and averaged. In XGBoost \cite{chen2016}, trees are
added one after another to reduce a regularised squared loss with second-order steps. A separate
model was trained for each crop, because the crops differ in season, yield level and weather
response. Hyperparameters were not tuned per condition, so that differences between conditions
could not be caused by tuning.

\begin{table}[!t]
\caption{Models and Fixed Hyperparameters}\label{tab:models}
\centering\footnotesize
\begin{tabular}{lp{5.6cm}}
\toprule
Model & Settings \\
\midrule
Mean & Mean of the training yields \\
Ridge & Standardised inputs; $\alpha \in \{10^{-3}, \ldots, 10^{3}\}$ (13 values) by internal CV \\
RF & 300 trees, min.\ 2 samples per leaf, 50\% of features per split \\
XGB & 400 trees of depth 4, learning rate 0.05, 80\% row and column subsampling, min.\ child weight 3 \\
\bottomrule
\end{tabular}
\end{table}

\subsection{Evaluation Designs}

Three designs were used. In the \emph{random split}, the observations of a crop were divided into
five folds at random, as in most earlier studies; districts and years seen in training therefore
also appear in the test folds. In the \emph{temporal split}, the model was trained on the seasons
1997--98 to 2015--16 and tested on 2016--17 to 2022--23 (96 to 205 test observations per crop),
which measures transfer to future years in familiar districts. In the \emph{district holdout},
the districts were divided into five groups, and each group was predicted by a model trained on the
other four, so that no observation of a test district was used for training, imputation or model
fitting. Every condition was repeated with three seeds, which changed the folds (except in the
temporal split, which is fixed), the missingness masks and the model randomness.

\subsection{Additional Experiments}

Two further experiments used the district holdout. In the \emph{feature-selection} experiment,
the full set of 18 features was compared with eight features chosen on each training fold by
mutual information \cite{kraskov2004} or by recursive feature elimination with a random forest
\cite{guyon2002}, at MCAR rates of 0, 20 and 40\%; selection was carried out on a median-imputed
copy of the training data. In the \emph{location} experiment, the latitude and longitude of the
district headquarters were added to the features of complete data.

\subsection{Implementation Details}

The experiments were implemented in Python 3.13 with scikit-learn 1.9.1 \cite{pedregosa2011},
XGBoost 3.4.1, NumPy 2.5, pandas 3.0 and SciPy 1.18. All runs were performed on a laptop CPU
(Intel Core i7-1255U, 12 threads, 16~GB memory) without a GPU, with conditions run in parallel.
Every random step was controlled by a seed, and all predictions were stored so that every reported
number can be recomputed from the released files.
